\section{Indicadores de desempe\~no (KPIs)}

Los KPIs se definieron para decidir qu\'e combinaci\'on de t\'ecnica,
\textit{view} y perfil conviene recomendar. La decisi\'on no depende solo del
retorno: tambi\'en considera riesgo, permanencia del cliente y utilidad para
FinPUC.

\begin{itemize}
    \item \textbf{Retorno total rebalanceado:}
    $R_{\text{total}} = \prod_t (1 + r_t) - 1$. Mide la ganancia acumulada
    fuera de muestra. Se usa para comparar el resultado final de cada cartera.

    \item \textbf{Retorno anualizado:}
    $R_{\text{anual}} = (1 + R_{\text{total}})^{1/T} - 1$. Normaliza el retorno
    por a\~no y permite comparar horizontes de distinta duraci\'on.

    \item \textbf{Volatilidad anualizada:}
    $\sigma_{\text{anual}} = \sigma_{\text{diaria}}\sqrt{252}$. Mide la
    variabilidad del retorno. Sirve para descartar alternativas demasiado
    inestables para el perfil del cliente.

    \item \textbf{Ratio de Sharpe rebalanceado:}
    $\text{Sharpe} = (R_{\text{anual}} - R_f)/\sigma_{\text{anual}}$, con
    $R_f = 0.02$. Resume retorno ajustado por riesgo. Es el principal KPI para
    comparar eficiencia financiera entre t\'ecnicas.

    \item \textbf{M\'aximo drawdown:} mayor ca\'ida porcentual desde un m\'aximo
    acumulado de riqueza. Eval\'ua la p\'erdida que el cliente habr\'ia tenido
    que tolerar, por lo que pesa especialmente en perfiles conservadores.

    \item \textbf{CVaR diario al 95\%:} p\'erdida promedio en el 5\% peor de
    d\'ias \cite{rockafellar2000}. Se usa para medir riesgo de cola y evitar
    carteras con p\'erdidas extremas demasiado severas.

    \item \textbf{Riqueza terminal esperada:} valor final promedio desde
    $C_0 = 1\,000$ USD en las 5\,000 simulaciones Monte Carlo. Traduce el
    desempe\~no financiero a una cifra monetaria para el cliente.

    \item \textbf{Probabilidad de ganancia:} proporci\'on de simulaciones en
    que la riqueza terminal supera $C_0$. Indica qu\'e tan probable es que el
    cliente termine el horizonte con utilidad.

    \item \textbf{Tasa de retiro:} proporci\'on de simulaciones en que el
    cliente abandona antes del final por superar su tolerancia de p\'erdida.
    Afecta la decisi\'on porque reduce la continuidad del cliente y las
    comisiones futuras.

    \item \textbf{Utilidad por comisiones:} ingreso esperado de FinPUC con
    comisi\'on $k = 1\%$, incluyendo comisi\'on inicial y comisiones por
    recomendaciones aceptadas. Mide el atractivo de la alternativa para la
    empresa.

    \item \textbf{Tasa de aceptaci\'on de recomendaciones:} proporci\'on de
    recomendaciones aceptadas seg\'un $P_2$. Una tasa baja indica que la cartera
    no supera el umbral de atractivo del perfil, reduciendo rotaci\'on y utilidad.
\end{itemize}

Para ordenar las alternativas se usa el Score P4:
\begin{equation}
\text{Score}_{P4} =
E[\text{Riqueza terminal}]
+ E[\text{Utilidad empresa}]
- C_0 \cdot \text{Tasa de retiro}.
\end{equation}

El sistema recomienda la combinaci\'on con mayor Score P4, siempre que sus
KPIs de riesgo sean consistentes con el perfil evaluado. As\'i, la decisi\'on
final equilibra valor para el cliente, riesgo asumido, retenci\'on y utilidad
para FinPUC.
